\chapter{Continuum Module Cleanup and the Finite-Strain Constitutive Interface}
\label{ch:continuum-constitutive-interface}

\section{Motivation}

The continuum module already contained the essential tensor infrastructure for
large displacements and finite strains:
\begin{itemize}
  \item deformation and stress tensors in \texttt{Tensor2}, \texttt{SymmetricTensor2},
        and \texttt{Tensor4};
  \item kinematic formulations in \texttt{SmallStrain},
        \texttt{TotalLagrangian}, and \texttt{UpdatedLagrangian};
  \item hyperelastic models written directly in tensor form.
\end{itemize}

However, at the \texttt{ContinuumElement} /\texttt{Material<>} boundary, the
contract was still essentially
\[
  \texttt{engineering Voigt strain} \;\longrightarrow\; \texttt{stress/tangent}.
\]
This was acceptable for infinitesimal strain and also acceptable as an adapter
for Green--Lagrange hyperelasticity, but semantically it was still too narrow:
the constitutive side could not explicitly see the deformation gradient
$\mathbf{F}$, the active stress measure, or the distinction between
infinitesimal and Green--Lagrange strain.

That is the exact point where future finite-strain plasticity, large
displacement RC continuum models, and eventually crack/fracture models will
need a cleaner interface.

\section{Problem Statement}

Before this refactor, the continuum path looked like:
\[
  \texttt{KinematicPolicy::evaluate()}
  \;\Rightarrow\;
  \texttt{gp.strain\_voigt}
  \;\Rightarrow\;
  \texttt{StateVariableT}
  \;\Rightarrow\;
  \texttt{MaterialPoint::compute\_response(StateVariableT)}.
\]

This had three architectural drawbacks:
\begin{enumerate}
  \item The element knew only a flat constitutive input, not the richer
        continuum kinematics from which it came.
  \item Finite-strain relations were forced to masquerade as small-strain
        \texttt{Strain<N>} relations even when the real primitive was
        $\mathbf{F}$ or $\mathbf{E}$.
  \item The interface did not clearly document which work-conjugate pair was
        active for each formulation.
\end{enumerate}

\section{New Carrier: \texttt{ConstitutiveKinematics}}

To fix that semantic compression, the continuum module now provides
\texttt{ConstitutiveKinematics<dim>}. It is a continuum-only carrier that
contains:
\begin{itemize}
  \item the active engineering strain vector used by the current formulation;
  \item the infinitesimal strain tensor;
  \item the Green--Lagrange strain tensor;
  \item the deformation gradient $\mathbf{F}$;
  \item $\det \mathbf{F}$;
  \item explicit measure tags:
        \texttt{StrainMeasureKind} and \texttt{StressMeasureKind}.
\end{itemize}

This carrier is not a replacement for \texttt{Strain<N>}. It is the
continuum-side bridge that preserves the old interface as a fallback while
opening a better path for finite kinematics.

\section{Mapping from Kinematic Policies}

The helper
\[
  \texttt{make\_constitutive\_kinematics<KinematicPolicy>(gp)}
\]
maps the Gauss-point bundle returned by \texttt{KinematicPolicy::evaluate()}
into the richer constitutive carrier.

\subsection*{Small strain}

For \texttt{SmallStrain}, the carrier is interpreted as:
\begin{align}
  \boldsymbol{\varepsilon} &= \nabla^s \mathbf{u}, \\
  \mathbf{F} &= \mathbf{I}, \\
  \det \mathbf{F} &= 1,
\end{align}
with active measure
\[
  (\text{infinitesimal strain},\; \text{Cauchy stress}).
\]

\subsection*{Total Lagrangian}

For \texttt{TotalLagrangian}, the carrier is interpreted as:
\begin{align}
  \mathbf{F} &= \mathbf{I} + \nabla_0 \mathbf{u}, \\
  \mathbf{E} &= \frac{1}{2}(\mathbf{F}^T\mathbf{F} - \mathbf{I}),
\end{align}
with active measure
\[
  (\text{Green--Lagrange strain},\; \text{2nd Piola--Kirchhoff stress}).
\]

\subsection*{Updated Lagrangian}

For \texttt{UpdatedLagrangian}, the constitutive carrier is now interpreted as:
\begin{align}
  \mathbf{e} &= \frac{1}{2}\left(\mathbf{I} - \mathbf{b}^{-1}\right), \\
  \mathbf{b} &= \mathbf{F}\mathbf{F}^T,
\end{align}
with active measure
\[
  (\text{Almansi strain},\; \text{Cauchy stress}).
\]

This is the spatial counterpart of the total-Lagrangian pair
$(\mathbf{E},\mathbf{S})$ and it is the natural route for current-configuration
constitutive laws.

\section{Optional Direct Continuum Interface}

The constitutive side now has an optional refinement:
\begin{quote}
\texttt{ContinuumKinematicsAwareConstitutiveRelation}
\end{quote}

If a constitutive relation implements
\begin{align}
  \texttt{compute\_response(ConstitutiveKinematics<dim>)} \\
  \texttt{tangent(ConstitutiveKinematics<dim>)}
\end{align}
then the erased \texttt{Material<>} handle will forward the richer carrier
directly. Otherwise it falls back to the legacy conversion
\[
  \texttt{ConstitutiveKinematics} \;\to\; \texttt{StateVariableT}.
\]

This keeps the existing code working while making the finite-strain interface
available immediately.

\section{Hyperelastic Relations}

\texttt{HyperelasticRelation<Model>} now implements the direct continuum
carrier overload. Internally it consumes the Green--Lagrange strain tensor
stored in the carrier:
\[
  \mathbf{S} = \frac{\partial W}{\partial \mathbf{E}}, \qquad
  \mathbb{C} = \frac{\partial^2 W}{\partial \mathbf{E}^2}.
\]

This change is architecturally important because it removes the need for a
finite-strain law to pretend that its natural interface is only a Voigt vector.
The old \texttt{Strain<N>} adapter is preserved, but it is no longer the only
semantic entry point.

\section{Changes in \texttt{ContinuumElement}}

\texttt{ContinuumElement} now evaluates kinematics as before, but it no longer
immediately collapses them into a flat constitutive input. The local assembly
loop now does:
\begin{enumerate}
  \item evaluate the Gauss-point kinematics through the selected
        \texttt{KinematicPolicy};
  \item build a \texttt{ConstitutiveKinematics<dim>} carrier;
  \item ask the type-erased material point for stress and tangent through that
        carrier;
  \item fall back automatically to \texttt{StateVariableT} if the constitutive
        law does not provide a richer overload.
\end{enumerate}

This is the correct direction for large-displacement and finite-strain
continuum modelling because the element is no longer hardwired to the idea that
``constitutive input = one Voigt vector''.

\section{Compile-Time and Performance Considerations}

The refactor was designed so that:
\begin{itemize}
  \item the hot path inside a concrete constitutive law remains statically
        dispatched;
  \item the erased boundary still pays the same category of cost as before
        (virtual call through \texttt{Material<>});
  \item legacy integrators that only know \texttt{KinematicT} remain valid;
  \item direct continuum overloads are purely additive.
\end{itemize}

In other words, the new interface broadens the semantic contract without
forcing a runtime penalty inside the concrete law implementation.

\section{Limitations After This Step}

This cleanup does \emph{not} mean the full finite-strain continuum stack is
finished.

The main remaining limitations are:
\begin{itemize}
  \item the spatial route is now explicit for continuum hyperelastic laws, but
        finite-strain inelastic spatial laws still need their own local-state
        and return-algorithm infrastructure;
  \item there is not yet an explicit continuum-specific
        \texttt{ConstitutiveState} split for finite-strain plasticity or damage;
  \item finite-strain internal-variable evolution still needs a more general
        residual/Jacobian-based local integration contract beyond the current
        small-strain J2-oriented path.
\end{itemize}

\section{Tests Added}

The new test suite verifies:
\begin{itemize}
  \item correct construction of \texttt{ConstitutiveKinematics} from
        \texttt{SmallStrain} and \texttt{TotalLagrangian};
  \item direct continuum-carrier consistency for
        \texttt{SVKRelation<3>} and \texttt{NeoHookeanRelation<3>};
  \item correct propagation of the new interface through the erased
        \texttt{Material<ThreeDimensionalMaterial>} handle.
\end{itemize}

\section{Architectural Outcome}

This step completes an important semantic cleanup of the continuum module:
\begin{quote}
kinematics, constitutive driving variables, and constitutive law entry points
are no longer implicitly collapsed into the same object.
\end{quote}

That is a necessary precondition for the next stages:
\begin{itemize}
  \item finite-strain plasticity with explicit internal state;
  \item large-displacement RC continuum models;
  \item damage/plasticity couplings;
  \item crack-band or other regularised fracture models for continuum RC.
\end{itemize}
